% Standalone section fragment.  Requires amsmath, amssymb, amsthm.
% The theorem concerns diagonal trace equations only.
\section{A positive Gaussian network for the diagonal trace equations}
\label{sec:gaussian-trace-network}

We work at the product point
\[
 A=E_i^8=E_i^4\times E_i^4,
 \qquad E_i=\mathbb C/(\mathbb Z+i\mathbb Z),
\]
with its product principal polarization \(\theta\).  The balanced
Gaussian action is multiplication by \(i\) on the first block and by
\(-i\) on the second.  Write \(U\) and \(V\) for these two sets of four
coordinates.  A divisor is represented by a Hermitian matrix, and
\(\theta\) by the identity matrix.  Put
\[
 D_0=\operatorname{diag}(1,-1,2,-2,3,-3,4,-4),\qquad
 F(t)=1-\frac{15}{28}t^2+\frac{13}{80}t^4
       -\frac{41}{1008}t^6+\frac1{70}t^8.
\]
For a unitary frame \(g\), its complete permutation packet consists
of \(gP_\sigma D_0P_\sigma^*g^*\), \(\sigma\in S_8\).
The elementary symmetric-polynomial identity gives
\begin{equation}
 \sum_{\sigma\in S_8}
   \exp(gP_\sigma D_0P_\sigma^*g^*)=8!F(\theta).
 \label{eq:trace-packet-average}
\end{equation}

\begin{theorem}[A positive network satisfying the diagonal equations]
\label{thm:gaussian-positive-trace-network}
There is a finite connected bipartite graph of distinct rational
trace-zero Hermitian divisors on \(A\), with a positive integer
multiplicity at every vertex, having the following properties.
Every edge has a rank-two difference of signature \((1,1)\), and
the weighted sum of exponentials is
\[
 R F(\theta)+c w,
 \qquad R>0,\quad c\ne0,\quad 0\ne w\in W_K(A).
\]
For every polarized Weil tangent direction \(\kappa\), the vector of
weighted diagonal obstructions admits an edge flow:
\begin{equation}
 \bigl(m_x\,\kappa\mathbin{\lrcorner}D_x\bigr)_x
 \in\operatorname{im}\left(
   \bigoplus_{e}q_e\longrightarrow\bigoplus_x H^2(A,\mathcal O_A)
 \right),
 \label{eq:positive-trace-flow}
\end{equation}
where \(q_e\) is the antiholomorphic surface wedge of the edge and
the map has opposite signs at its endpoints.

After a common integral scaling, the vertices are represented by
line bundles.  Their positive multiplicities admit a single ordinary
extension that is locally free and simple, with Chern character
\(R F(N\theta)+cN^4w\).  The edge flows cancel the diagonal trace
equations for that extension.  With a suitable matrix-unit choice of
its extension entries they also cancel all matrix-valued diagonal
cohomology blocks.

No assertion is made that the full Atiyah obstruction vanishes, that
the bundle deforms in the Weil family, or that its semiregularity map
is injective.  The off-vertex equations and the strong
semiregularity condition remain unresolved.
\end{theorem}

\subsection{The finite group and its character}

Let \(G(4,4,4)\) be the group of monomial unitary \(4\times4\)
matrices with nonzero entries in \(\mu_4\) whose product is one.
Set
\[
 H=G(4,4,4)_U\times G(4,4,4)_V,
 \qquad \chi(h)=\det(h_U)\det(h_V)\in\{1,-1\}.
\]
The order of each factor is \(4^3\cdot4!=1536\).
The central scalar subgroup acting trivially on divisor matrices is
\(Z=\mu_4 I_8\), and \(\chi|_Z=1\).

\begin{lemma}\label{lem:gaussian-trace-character}
The \(\chi\)-isotypic subspace of
\(\bigoplus_p H^{p,p}(A)\) is exactly \(W_K(A)_\mathbb C\), in
degree \((4,4)\).
\end{lemma}

\begin{proof}
In one block a cohomology monomial has phase exponents in
\(\{-1,0,1\}^4\).  Invariance under the product-one diagonal
\(\mu_4\)-subgroup forces the four exponents to be equal modulo
four.  They are therefore all \(+1\), all \(-1\), or all zero.
The first two possibilities are the holomorphic and antiholomorphic
volume forms, each transforming by the determinant sign.
In the zero case, each coordinate occurs in both Hodge factors or
in neither.  The residual permutation representation is a module on
\(k\)-subsets of four coordinates.  It has no sign constituent:
the stabilizer \(S_k\times S_{4-k}\) contains an odd permutation
for every \(k\).  Thus only the two volume forms survive in each
block.  Equal holomorphic and antiholomorphic degree requires
opposite choices in the two blocks, giving exactly the two Weil
summands.
\end{proof}

There are twenty-four rank-one unitary involutions in each factor:
choose a coordinate pair and a phase in \(\mu_4\), and interchange
the two coordinates with inverse phases.  All forty-eight
involutions have \(\chi=-1\), and they generate \(H\): ordinary
interchanges generate the permutation subgroups, and products of
two phase interchanges generate the product-one diagonal subgroups.

\subsection{The exact frame and its packet stabilizer}

Use
\[
 g=h_{\mathbf1}h_{(1,2,3,4,5,6,7,8)},\qquad
 h_u=I-\frac{2uu^t}{u^tu}.
\]
Equivalently \(g=G/408\), where
\begin{equation}
 G_{ij}=408\delta_{ij}-102+(36-4i)j,
 \qquad 1\le i,j\le8.
 \label{eq:gaussian-exact-frame}
\end{equation}
The identity \(G^tG=408^2I\) is exact.  Every entry of \(G\)
is nonzero: an off-diagonal entry is \(2\) modulo \(4\), and the
diagonal entries are positive.

\begin{lemma}[No unintended packet identifications]\label{lem:gaussian-trace-frame}
The stabilizer in \(H\) of the unordered eigenline frame \(g\)
is exactly \(Z\).  Consequently the divisor vertices
\[
 D_{h,\sigma}=hgP_\sigma D_0P_\sigma^*g^*h^*,
 \qquad (hZ,\sigma)\in(H/Z)\times S_8,
\]
are distinct, and the coloring \(\chi(h)(-1)^\sigma\) is
well defined.
\end{lemma}

\begin{proof}
Suppose a monomial unitary matrix \(h\) permutes the eigenlines of
\(g\).  The columns of \(G\) have the same Euclidean norm, so the
proportionality factors induced by \(h\) have absolute value one.
Monomial unitaries preserve the sum of absolute coordinate values.
For the eight columns of \(G\), these sums are respectively
\[
 940,\ 844,\ 732,\ 676,\ 688,\ 732,\ 784,\ 832.
\]
The first column is the unique maximum, hence its line is fixed.
Its eight absolute coordinate values are
\[
 338,\ 74,\ 78,\ 82,\ 86,\ 90,\ 94,\ 98,
\]
which are pairwise distinct.  The row permutation of \(h\) is
therefore the identity.  Since all eight coordinates are nonzero,
fixing this line forces all diagonal phases to be equal.  Thus
\(h\) is scalar.  For \(h\in H\), the scalar belongs to \(\mu_4\),
and conversely every such scalar lies in \(H\) and preserves the
frame.  This proves the assertion without an enumeration.
\end{proof}

Retain all transposition edges within each packet and all full
matchings \((hZ,\sigma)\leftrightarrow(hrZ,\sigma)\) for the
forty-eight reflections \(r\).  Every matching is induced by a
rank-one unitary conjugation and therefore has a rank-two,
trace-zero Hermitian difference.  It is nonzero by the stabilizer
lemma, so its signature is \((1,1)\).  Internal transpositions and
external matchings both reverse the displayed coloring.  The graph
is connected because the two sets of generators connect the two
indices.

\subsection{A finite certificate for the reduced trace equations}
\label{subsec:gaussian-trace-certificate}

We use the convention that a Hermitian divisor is conjugated as
\(D\mapsto hDh^*\), while a polarized Kodaira--Spencer class is represented
by a complex symmetric matrix \(B\).  In a unitary eigenbasis \(g\), its
matrix is \(g^tBg\).  For a physical root normal \(u\), put
\[
 v=g^t\overline u.
\]
The antiholomorphic bridge line has, in this eigenbasis, the coordinates
\[
 q_{u,\sigma;kl}
   =v_kv_l(\lambda_{\sigma,l}-\lambda_{\sigma,k}).
\]
The sign and a common nonzero normalization of \(q\) are immaterial,
since its coefficient is a free flow variable.

Let \(\mathcal P\) be the primitive diagonal fourth-degree representation.
We realize it in the space with basis \(x_I\), \(|I|=4\), where
\(x_i^2=0\), as the fourteen-dimensional span of
\[
 p_M=\prod_{\{a,b\}\in M}(x_a-x_b)
\]
for the noncrossing perfect matchings \(M\) of
\(\{1,\ldots,8\}\).  Order these matchings recursively by pairing the
first unused index with each later index of opposite parity, and then
ordering the two remaining intervals recursively.  Let \(C\) be the
resulting \(70\) by \(14\) integer matrix.  Its first column is
\[
 p_0=(x_1-x_2)(x_3-x_4)(x_5-x_6)(x_7-x_8).
\]
The matrix \(R_s\) of a coordinate transposition \(s=(kl)\) is uniquely
determined by
\[
 CR_s=sC.
\]
All these matrices are integral.  Write \(D_s=I-R_s\); each \(D_s\)
has rank five.

We use the exact frame \eqref{eq:gaussian-exact-frame} throughout.
The identity \(G^tG=408^2I\) is included in the verification.

The roots are indexed by a pair \(a<b\) in one of the two blocks
\(\{1,\ldots,4\}\), \(\{5,\ldots,8\}\), and by
\(\zeta\in\{1,-1,i,-i\}\):
\begin{equation}
 u_{ab}(\zeta)=e_a-\zeta^{-1}e_b,\qquad
 r_{ab}(\zeta)=I-u_{ab}(\zeta)u_{ab}(\zeta)^*.
 \label{eq:gaussian-trace-roots}
\end{equation}
Thus \(r_{ab}(\zeta)\) is a rank-one unitary involution with determinant
\(-1\).  The normal has squared length two, explaining the coefficient
in \eqref{eq:gaussian-trace-roots}.

For a root in the first block, its allowed coefficient map is written
\begin{equation}
 a_u(B)=\sum_{b=5}^{8}(u^tBe_b)\,z_{u,b},
 \qquad z_{u,b}\in\mathcal P.
 \label{eq:gaussian-trace-potential-u}
\end{equation}
For a root in the second block, it is
\begin{equation}
 a_u(B)=\sum_{a=1}^{4}(e_a^tBu)\,z_{u,a}.
 \label{eq:gaussian-trace-potential-v}
\end{equation}
These are precisely the linear maps satisfying
\(a_u(r_u^tBr_u)=-a_u(B)\).  This sign is required by the determinant
character and the reversal of a bridge.

For \(f=e_1\in\mathbb Q^{14}\), the finite system to be solved is
\begin{equation}
 D_s\left[
    (g^tBg)_{kl}f-
    \sum_u v_kv_l\,a_u(B)
 \right]=0,
 \qquad s=(kl).
 \label{eq:gaussian-trace-reduced-system}
\end{equation}
It suffices to impose \eqref{eq:gaussian-trace-reduced-system} for the sixteen matrices
\[
 B_{ab}=E_{ab}+E_{ba},\qquad
 1\le a\le4< b\le8.
\]
There are \(16\cdot28\cdot14=6272\) displayed scalar equations.
Selecting five independent rows from each \(D_s\) gives \(2240\)
retained row positions.  The flow has
\(48\cdot4\cdot14=2688\) coordinates.

For completeness, the Gaussian coefficients are removed from the
linear system as follows.  For each root pair, let \(F_\zeta\) denote
the corresponding coefficient-column block in \eqref{eq:gaussian-trace-reduced-system}.  Replace
\[
 F_1,\ F_{-1},\ F_i,\ F_{-i}
 \quad\text{by}\quad
 F_1,\ F_{-1},\ \operatorname{Re}F_i,\ \operatorname{Im}F_i.
\]
The latter matrix is rational and has the same complex column span.
If \(x,y\) are its real and imaginary coefficient vectors, the
original two root potentials are recovered as
\begin{equation}
 z_i=(x-iy)/2,\qquad z_{-i}=(x+iy)/2.
 \label{eq:gaussian-trace-real-imag}
\end{equation}
Multiplication by \(408^2\) makes every coefficient integral.

The calculation can be regenerated deterministically from these
finite inputs.  Order the sixteen pairs \((a,b)\), the twenty-eight
pairs \((k,l)\), and the four-subsets lexicographically.  Order the
flow coordinates first by the block \(U,V\), then the root pair
lexicographically, then by the four column types
\(F_1,F_{-1},\operatorname{Re}F_i,\operatorname{Im}F_i\), then by
the complementary coordinate in increasing order, and finally by
the fourteen recursively ordered matchings.  Over
\(\mathbb F_{1009}\), scan the columns of each \(D_s\) from left
to right, use the first available nonzero row as the next pivot,
and retain the five original row indices so obtained.  Stack the
resulting equations in the stated tangent and transposition order.
For the resulting augmented matrix \([A\mid b]\), process columns
in consecutive panels of size \(48\).  In each remaining panel,
ordinary forward elimination scans columns from left to right and
chooses the first available nonzero row.  Move the selected original
rows, in pivot order, above the unselected rows, preserving the
latter's order.  If \(P\) is the selected square pivot block,
multiply its rows by \(P^{-1}\) and eliminate its pivot columns
from all remaining rows before processing the next panel.
This produces \(2224\) pivot rows and columns, all in \(A\).
Set all other flow coordinates to zero, solve the selected square
system over \(\mathbb Q\), clear denominators, and test every
one of the original \(6272\) equations.  These instructions specify
the rational certificate without requiring its numerators as an
additional mathematical input.

The resulting certificate supplies a positive integer \(d\) and an integer
vector \(n\), supported on \(2224\) of the \(2688\) flow coordinates,
such that
\begin{equation}
 A n=d\,b.
 \label{eq:gaussian-trace-integer-certificate}
\end{equation}
Here \(A\) is the completely specified integer matrix of \eqref{eq:gaussian-trace-reduced-system}, after
the real--imaginary replacement, and \(b\) is its right-hand side for
\(f=e_1\).  The denominator \(d\) has \(5500\) bits.  A pivot minor
was selected over \(\mathbb F_{1009}\), a square rational system was
solved by Dixon lifting, and \eqref{eq:gaussian-trace-integer-certificate} was subsequently verified by exact
integer multiplication.  Thus neither rank equality modulo a prime
nor a floating-point residual is being used as the final inference.

The accompanying computational workbench supplies the portable
program \emph{Gaussian trace network verification} and its plain
integer certificate.  The program reconstructs \(C\),
all \(R_s\), \(G\), all root coefficients, and both sides of \eqref{eq:gaussian-trace-integer-certificate}
from the formulas above, using only Python's standard library.
It checks all \(6272\) displayed coordinates, not merely the
\(2240\) row positions used to select the rational solve.  The
certificate consists of the common denominator, the integer
numerators, and their coordinate indices.  A separate finite
packet-stabilizer enumeration is retained in the workbench as
an additional check of the elementary proof above.

The certificate also checks that the chosen primitive direction has
nonzero Weil projection.  With \(U=\{1,2,3,4\}\) and
\(V=\{5,6,7,8\}\), its complementary minor coefficient is
\begin{equation}
 \sum_{|I|=4}C_{I,1}\det(g_{U,I})\det(g_{V,I})
 =
 \frac{49}{41616}\ne0.
 \label{eq:gaussian-trace-minor}
\end{equation}
Before reduction the denominator in \eqref{eq:gaussian-trace-minor} is \(408^8\).  Its
nonvanishing is checked by integer determinant expansion.

Equations \eqref{eq:gaussian-trace-reduced-system}--\eqref{eq:gaussian-trace-minor} certify the reduced diagonal trace-incidence
conditions.  The reconstruction and their precise deformation-theoretic
limits are given next; no full deformation or semiregularity assertion
is inferred from this finite calculation.


\subsection{From the reduced solution to positive multiplicities}

We now explain why the finite linear system supplies a flow on the
actual graph, rather than only an abstract representation.
Let \(\mathcal P\) and \(C\) be as in the computation.  If
\(p_I=\prod_{i\in I}\lambda_i\), let \(v\in\mathcal P\)
be the primitive projection of the diagonal fourth-degree vector
\((p_I)\), and define the rational functional
\[
 \ell(f)=\sum_{|I|=4}p_I(Cf)_I.
\]
Then \(\ell(v)>0\).  For \(f\in\mathcal P\), put
\begin{equation}
 m_f(\sigma)=\ell(R_{\sigma^{-1}}f).
 \label{eq:trace-matrix-coefficient}
\end{equation}
The convention is chosen so that
\[
 m_f(\sigma)-m_f(s\sigma)=m_{(I-R_s)f}(\sigma).
\]
The packet character associated with these multiplicities is
\begin{equation}
 \sum_{\sigma\in S_8}m_f(\sigma)e^{P_\sigma D_0P_\sigma^*}
 =\frac{8!}{14}\ell(v)\,Cf,
 \label{eq:trace-schur-lift}
\end{equation}
as a pure total cohomology vector in degree four.
Indeed, the Specht module \((4,4)\) occurs only in degree four
among the subset permutation modules.  The remaining operator in
that irreducible module is scalar by Schur's lemma; its trace is
\(8!\ell(v)\), giving \eqref{eq:trace-schur-lift}.

For the particular vector \(f=e_1\), this is completely integral:
\begin{equation}
 m_f(\sigma)=\prod_{j=1}^4
  \bigl(\lambda_{\sigma^{-1}(2j-1)}
       -\lambda_{\sigma^{-1}(2j)}\bigr).
 \label{eq:trace-explicit-multiplicity}
\end{equation}
Its absolute value is at most \(625\).  In fact, the sum of the
four absolute differences is at most
\((4+3+2+1)-(-1-2-3-4)=20\), and the arithmetic--geometric
mean inequality gives \(5^4\).  Therefore
\begin{equation}
 m_{h,\sigma}=626+\chi(h)m_f(\sigma)
 \label{eq:trace-positive-multiplicity}
\end{equation}
is an integer between \(1\) and \(1251\) at every vertex.

For clarity, the Schur scalar in this normalization is
\[
 \frac1{16}\sum_{\sigma\in S_8}m_f(\sigma)^2
 =160225344.
\]
Here \(16\) is the squared norm of the sixteen-term vector
\(Ce_1\).  The square sum is obtained simply by evaluating
\eqref{eq:trace-explicit-multiplicity} on the \(8!\) permutations.
The exact complementary-minor coefficient of \(g(Ce_1)\) is
\(49/41616\).  Thus, if \(w\) denotes the normalized real Weil
class whose two complementary matrix entries are one, the weighted
character is
\begin{equation}
 \sum_{hZ,\sigma}m_{h,\sigma}e^{D_{h,\sigma}}
 =R F(\theta)+
  160225344\,|H/Z|\,\frac{49}{41616}\,w,
 \quad R=626\,|H/Z|\,8!.
 \label{eq:trace-positive-character}
\end{equation}
The coefficient is nonzero.  This proves the character assertion
without identifying a nonzero primitive class with a Weil class
merely by its degree: the finite-group character lemma does the
identification.

\subsection{Reconstruction of the edge flows}

Let \(z_r(\kappa)\in\mathcal P\) be the certified coefficient
map for a root reflection.  Its anti-invariance is
\[
 z_r(r^tBr)=-z_r(B),
\]
with \(B\) the symmetric matrix of \(\kappa\).  The inverse
transport of a tangent by \(h\) is \(B\mapsto h^tBh\).
At the edge joining \((hZ,\sigma)\) to \((hrZ,\sigma)\),
use the scalar potential
\[
 Z_{h,r,\sigma}(B)=
  \chi(h)\,m_{z_r(h^tBh)}(\sigma).
\]
This gives a well-defined flow on the quotient by \(Z\): replacing
\(h\) by \(hz\), \(z\in\mu_4\), multiplies the inverse-transported
tangent and hence the scalar potential by \(z^2\), and multiplies
the transported antiholomorphic wedge generator by \(\overline z^2\).
Their product is unchanged.
Its natural surface wedge at the second endpoint is the negative
of that at the first.  Since \(\chi(r)=-1\) and the coefficient
map is anti-invariant, the scalar potential itself is unchanged
when the endpoints are interchanged.  Thus the two contributions
are opposite, as required for an edge flow.

After subtracting all such external flows, the coefficient at a
vertex in the coordinate wedge \((k,l)\) is the eigenvalue
difference \(\lambda_{\sigma,l}-\lambda_{\sigma,k}\) times
\[
 (g^tBg)_{kl}m_f(\sigma)
   -\sum_r v_{r,k}v_{r,l}m_{z_r(B)}(\sigma).
\]
The certified equation says exactly that the bracket is invariant
under the left transposition \((kl)\).  The eigenvalue difference
changes sign at the other endpoint of the internal matching.
Hence this residual is canceled by that internal edge.  The
constant background \(626\) has the same elementary internal
matching flow.  This proves \eqref{eq:positive-trace-flow} for all
sixteen Weil tangent directions and hence for their span.

\subsection{Line bundles, simplicity, and the remaining obstruction}

Choose a common positive integer \(N\) making every Riemann form
\(ND_{h,\sigma}\) even integral.  One may first take
\(N_0=2\cdot408^2\): the remaining matrices are Gaussian
monomial conjugates of \(G P_\sigma D_0P_\sigma^tG^t\).
Use Appell--Humbert representatives with trivial semicharacter.
Their ratios on an edge descend to the actual abelian surface
quotient and have index one, by Mumford's Appell--Humbert theorem
and complex index theorem \cite[\S2, p.~20; \S16, pp.~150,162]{MumfordAV}.
Both edge Ext\({}^1\) groups are nonzero and pair perfectly into
the line \(q_e\).  Multiplying \(N_0\) by \(1251\) makes their
dimensions at least every product of endpoint multiplicities,
because the index-one dimension grows quadratically and is a
positive integer at \(N_0\).

Let \(T\) and \(B\) be the sums of the line bundles with their
multiplicities in the two color classes.  Distinct vertices have
no Hom: a nonzero divisor of polarization degree zero cannot be
effective.  For every oriented edge choose independent Ext\({}^1\)
classes \(\eta_{ab}\), one for each rectangular matrix unit, and
give the ordinary extension
\[
 0\longrightarrow B\longrightarrow\mathcal E\longrightarrow T
 \longrightarrow0
\]
the component \(\sum_{a,b}\eta_{ab}\otimes E_{ab}\).
This extension is locally free.  An endomorphism preserves \(B\)
because \(\operatorname{Hom}(B,T)=0\), and induces block matrices
satisfying \(b_vE_{ab}=E_{ab}t_u\) at every edge.  They must be
common scalar matrices.  Connectivity therefore gives
\(\operatorname{End}(\mathcal E)=\mathbb C\).

There is a useful strengthening of the scalar trace statement for
this chosen extension.  The two endpoint products of degree-one
classes have opposite traces by the graded trace identity.  After
incorporating the incidence minus sign at the second endpoint,
the reverse cup product on one edge maps onto the pairs of endpoint
matrices having equal trace.  Indeed, the dual kernel of this
sign-adjusted map is the space of pairs satisfying
\(A_vE_{ab}=E_{ab}A_u\) for every matrix unit, which consists of
common scalars.  For a certified oriented edge flow \(a q_e\),
prescribe the actual endpoint matrices
\(a q_e I_{m_v}/m_v\) and \(-a q_e I_{m_u}/m_u\).
Their sign-adjusted traces agree, so they are realizable by a
closed reverse Ext\({}^1\) class.  Summing over edges cancels
all matrix-valued diagonal cohomology blocks; exact diagonal forms
can then be removed by a diagonal Dolbeault correction.

This still does not solve the full deformation equation.  Off-vertex blocks within either color class can be created by
products along two different incident edges.  The upper block
between the two color classes also contains the differentiated
extension class and, in a minimal model, the corresponding
\(m_3(e,x,e)\) term.  Its vanishing has not been proved.  Nor has
\(\dim\operatorname{Ext}^2(\mathcal E,\mathcal E)=104\), or an
appropriate equivariant version, been established.  The rank
\(104\) of the Hochschild action on the character is the
\(n=4\), Hankel-rank-three specialization of the polarized
criterion of Bhattacharjee--Bhattacharya, Theorem~A6.  The finite
group \(H\) above is a construction symmetry at the CM point;
it is not being asserted to be a translation/Picard group extending
over the Weil family.  The remaining off-vertex deformation equations
and the semiregularity conditions require further arguments.
